% 出席確認クイズ 第01回　ガイダンス／AI時代の経営情報
%   attendance/attendance01.html から生成（解説は本ガイド用に加筆）。

\quizq{1}{「経営情報」という科目で扱う「情報」の位置づけとして最も適切なものはどれでしょうか?}%
  {社内メールの総称}%
  {\correct{ヒト・モノ・カネと並ぶ経営資源のひとつ}}%
  {会計帳簿の別名}%
  {通信回線の速度のこと}%
  {正解は (b)。経営資源はヒト・モノ・カネに「情報」を加えた4つで整理されることが多く、本講義はこの第4の資源としての情報を扱います。(a)(c)(d) はいずれも情報の一側面や伝達手段にすぎず、経営資源としての位置づけを表していません。情報が資源であるなら、集める・整える・活用する・守るという「経営」の対象になる、という発想が出発点です。}

\quizq{2}{生成AI(Generative AI)の説明として適切なものはどれでしょうか?}%
  {家電を遠隔操作するための技術}%
  {与えられたデータを分類するだけのAI}%
  {コンピュータウイルスを検知するソフト}%
  {\correct{文章・画像・音声などの新しいコンテンツを生成するAI}}%
  {正解は (d)。生成AIは、学習したパターンをもとに文章・画像・音声などを新たに「生成」する点が特徴です。(b) の分類や判別は従来型（識別系）AIの典型で、生成AIとは区別されます。(a)(c) はAIの応用先や別分野の技術であって、定義ではありません。第3回でしくみを詳しく扱います。}

\quizq{3}{AIの出力を経営判断に使うときの基本姿勢として適切なものはどれでしょうか?}%
  {AIの出力は常に正しいので、そのまま採用する}%
  {\correct{出力の根拠と限界を確認したうえで、人が判断する}}%
  {回答が長いほど信頼できると考える}%
  {AIは使わないのがもっとも安全である}%
  {正解は (b)。AI Policy の「検証する」にあたる姿勢です。(a) はハルシネーション（第3回）の危険を無視しています。(c) は回答の長さと正確さに関係がないため誤りです。(d) は「使う」という本講義の方針に反し、AIを使わないこと自体がリスク（競争力の低下、シャドーAIの発生）になる場合もあります。}

\quizq{4}{「AI」の正式名称はどれでしょうか?}%
  {\correct{Artificial Intelligence}}%
  {Automatic Information}%
  {Advanced Internet}%
  {Analog Interface}%
  {正解は (a)。AI は Artificial Intelligence（人工知能）の略です。他の選択肢は実在しない語の組み合わせです。語源を押さえると、AIは「人工の」知能であり、万能でも魔法でもなく、設計と学習データに依存した道具だと意識しやすくなります。}

\quizq{5}{本講義の「AI Policy(AIとの向き合い方)」として掲げられて\underline{いない}ものはどれでしょうか?}%
  {AIの寄与と自分の判断を区別して示す}%
  {毎回の演習でAIを実際に使う}%
  {\correct{AIの利用を隠す}}%
  {出力を検証してから採用する}%
  {正解は (c)。本講義の AI Policy は「使う・検証する・区別する」の3つです。AI の利用は隠すのではなく、AI活用記録で明示します。(b) は「使う」、(d) は「検証する」、(a) は「区別する」にそれぞれ対応します。利用の透明性は評価の対象でもあります。}
